\section{ERA: evidence-guided evaluator evolution}
\label{sec:method}

\subsection{An executable evaluator}
\label{sec:runtime}

ERA changes how a program obtains and uses evidence. Artifact adapters produce
representations appropriate to the domain, such as rendered views, source text,
spatial structure, or interaction traces. Routing selects tools and sends their
outputs to evaluator modules with explicit responsibilities. The modules apply
the rubric and return findings with provenance and uncertainty. Fusion combines
these findings into an overall score and critique
(Figure~\ref{fig:runtime}). Existing metrics are available as tools when their
input contracts are satisfied; availability does not require every tool to run
on every artifact.

\begin{figure}[t]
\centering
\begingroup
\renewcommand{\sfdefault}{phv}% Match the vector charts without changing body fonts.
\begin{tikzpicture}[line cap=round,line join=round]
  \path[use as bounding box] (0,0) rectangle (13.7,4.95);
  \pPanel{a}{Read one artifact}{0}{4.72}
  \pPanel{b}{Route observations}{3.37}{4.72}
  \pPanel{c}{Apply criteria}{7.01}{4.72}
  \pPanel{d}{Synthesize}{10.62}{4.72}
  \path[pSurface,fill=pBlueWash] (0,1.02) rectangle (2.93,4.38);
  \path[pSurface,fill=pTealWash] (3.81,1.02) rectangle (6.65,4.38);
  \path[pSurface,fill=pGoldWash] (7.08,1.02) rectangle (10.17,4.38);
  \path[pSurface,fill=pEvoWash] (10.71,1.02) rectangle (13.70,4.38);

  \pIcon{page}{pBlue}{.34}{4.03}{.58}
  \node[pSmall,text=pBlue,anchor=west] at (.70,4.07) {$A$: artifact views};
  \path[draw=pBlue!30,fill=white,rounded corners=3pt]
    (.17,1.12) rectangle (2.76,3.61);
  \node[pLabel,text=pBlue,anchor=west] at (.28,3.39) {Summary $y$};
  \node[pText,anchor=north west] at (.28,3.15)
    {Admissions fell for\\\textcolor{pRed}{\underline{all age groups.}}};
  \draw[pRule] (.31,2.32) -- (2.60,2.32);
  \node[pLabel,anchor=west] at (.28,2.09) {Task source $x$};
  \node[pSmall,text=pInk,anchor=north west] at (.28,1.85)
    {Adult admissions fell.\\Children not studied.};

  \pIcon{loupe}{pTeal}{4.12}{4.03}{.59}
  \node[pSmall,text=pTeal,anchor=west] at (4.46,4.07) {$T$: permitted tools};
  \node[circle,draw=pTeal!55,fill=white,minimum size=.56cm,
    inner sep=0pt,pTitle,text=pTeal] (route) at (3.38,2.70) {$\Pi$};
  \draw[pArrow,draw=pBlue] (2.77,2.70) -- (route.west);
  \foreach \y/\id/\call/\value in {
    3.34/scope/{scope(y)}/{all age groups},
    2.38/source/{source(x)}/{adults only},
    1.43/raw/{text(y)}/{original wording}}{
    \node[pCard,draw=pTeal!28,fill=white,text width=2.18cm,
      minimum height=.73cm,align=left] (\id) at (5.25,\y)
      {\textcolor{pTeal}{\texttt{\call}}\\\value};
    \draw[pEvidence] (route.east) to[out=0,in=180] (\id.west);
  }

  \pIcon{clipboard}{pGold}{7.41}{4.03}{.59}
  \node[pSmall,text=pGold,anchor=west] at (7.76,4.07) {$K$: domain criteria};
  \path[draw=pGold!25,fill=white,rounded corners=4pt]
    (7.23,2.53) rectangle (10.01,3.65);
  \node[pLabel,text=pGold,anchor=west] at (7.36,3.44) {$V_1$\quad Faithfulness};
  \node[pText,anchor=west] at (7.36,3.07) {Scope overreach};
  \node[pSmall,anchor=west] at (7.36,2.73) {Source $x$; $y$, clause 2};
  \draw[pEvidence] (scope.east) to[out=0,in=180] (7.22,3.20);
  \draw[pEvidence] (source.east) to[out=0,in=180] (7.22,2.80);
  \path[draw=pGold!25,fill=white,rounded corners=4pt]
    (7.23,1.18) rectangle (10.01,2.18);
  \node[pLabel,anchor=west] at (7.36,1.97) {$V_2$\quad Fluency};
  \node[pText,anchor=west] at (7.36,1.61) {Readable sentence};
  \node[pSmall,anchor=west] at (7.36,1.32) {At $y$, sentence 1};
  \draw[pEvidence] (raw.east) to[out=0,in=180] (7.22,1.57);

  \node[pSmall,align=center] at (12.20,4.06) {Evidence $\ne$ judgment};
  \draw[pEvolution] (10.01,3.07) to[out=0,in=155] (10.96,2.69);
  \draw[pEvolution] (10.01,1.61) to[out=0,in=205] (10.96,2.61);
  \node[circle,draw=pEvo!55,fill=white,minimum size=.53cm,
    inner sep=0pt,pTitle,text=pEvo] (fuse) at (11.18,2.65) {$G$};
  \pIcon{clipboard}{pEvo}{12.64}{3.29}{.84}
  \node[pText,align=center] at (12.64,2.62) {$s_\theta(x,y)$};
  \node[pLabel,align=center] at (12.64,2.26) {Score + critique};
  \draw[pEvolution] (fuse.east) -- (11.94,2.65);
  \node[pSmall,align=center] at (12.20,1.64) {Findings + trace $\tau_\theta$};
  \node[pSmall,text=pEvo,align=center] at (12.20,1.22) {Numerical / semantic $G$};

  \draw[pRule] (0,.76) -- (13.7,.76);
  \node[pText,align=center] at (6.85,.48)
    {\textbf{Editable:}\quad $A$ views \quad $K$ rubrics \quad $V$ modules \quad
     $T$ tools \quad $\Pi$ routing \quad $G$ fusion};
  \node[pSmall,align=center] at (6.85,.12)
    {Fixed model; no opponent, human labels, mining labels, split identity, or generator identity at runtime};
\end{tikzpicture}
\endgroup

\caption{The six-component evaluation program. (a) Read one candidate and
permitted task evidence. (b) Route tool observations. (c) Apply dimensional
criteria. (d) Fuse sourced findings into a score and critique. The lower rail
identifies the editable components. Human feedback supervises evolution, never
the deployed scorer. Evidence and judgment remain separate; the summary and
findings are a constructed example.}
\label{fig:runtime}
\end{figure}

Fusion may use a declared numerical rule or semantic synthesis. Semantic synthesis
reads specialist findings and their anchored quality scales, groups overlapping
or conflicting factors, and locates overall quality on its own scale. It does
not average raw scores with different meanings. References to findings must be
accounted for, although correct reference bookkeeping does not establish semantic
correctness. Appendix~\ref{app:synthesis} illustrates this distinction.

\subsection{Direction hypotheses and bounded revisions}
\label{sec:proposal}

An adaptation direction states a mechanism that could explain a discrepancy with
human judgment. For example, a summary evaluator may miss lost qualifications
because it checks individual facts without comparing their scope against the
source. A direction specifies the relevant evidence, competing explanations,
target and protection cases, editable components, and a falsifiable prediction.
The proposal process examines original training artifacts and repeated execution
traces; compact summaries are navigation aids rather than substitutes for evidence.

Each elementary edit changes an attributable component of $(A,K,V,T,\Pi,G)$.
When a mechanism requires a new tool, a route, and a consumer change, the edits
may be assembled into one connected candidate with a fixed parent. The complete
candidate is measured once as a system; its outcome is not credited separately
to every edit. Edit-size and component constraints bound the search action,
but do not guarantee a small semantic change. Evaluator modules may be split,
merged, or removed as their responsibilities change.

\subsection{Proposal, measurement, and allocation}
\label{sec:search}

ERA separates proposal evidence from candidate selection. The proposal subset
$\Dprop$ exposes authorized training artifacts, available human feedback, and
versioned traces. A source-disjoint, sealed training subset $\Ddev$ supplies
aggregate development measurements only. Local checks on $\Dprop$ assess
predictions and regressions without estimating population performance.
A frozen candidate is also measured on the common $\Ddev$ cohort.
Validation $\Dval$ is terminal; Test/OOD remain locked until the method and
analysis are frozen (Appendix~\ref{app:boundaries}).

\begin{figure}[t]
\centering
\begin{tikzpicture}[line cap=round,line join=round]
  \path[use as bounding box] (0,0) rectangle (13.7,4.80);
  \pPanel{a}{One hypothesis, connected edits}{0}{4.57}
  \pPanel{b}{Mechanism checks}{6.46}{4.57}
  \pPanel{c}{Allocate or stop}{10.33}{4.57}
  \node[pText,anchor=west] at (.02,3.98)
    {Hypothesis: the evaluator lacks scope evidence.};

  \node[pDot,draw=pInk,fill=pInk,text=white,minimum size=.58cm] (b) at (.40,2.80) {$B$};
  \node[pSmall,align=center] at (.40,3.39) {Fixed\\parent};
  \draw[pPlum,dashed,rounded corners=2pt] (1.19,2.02) rectangle (6.06,3.47);
  \foreach \x/\id/\s/\t in {1.90/t/T/{Add scope\\extractor},
    3.57/pi/\Pi/{Route the\\source span},5.23/v/V/{Compare\\qualifiers}}{
    \node[draw=pPlum,fill=pPlum!7,minimum size=.50cm,inner sep=0pt,pTitle,text=pPlum]
      (\id) at (\x,2.96) {$\s$};
    \node[pSmall,align=center,text=pInk] at (\x,2.36) {\t};
  }
  \draw[pRevision] (b.east) -- (1.04,2.80) -- (1.04,2.96) -- (t.west);
  \draw[pRevision] (t) -- (pi);
  \draw[pRevision] (pi) -- (v);
  \node[pSmall,text=pPlum,fill=white,inner sep=2pt] at (3.59,3.47)
    {$W=\operatorname{patch}(B)$};
  \node[pDot,draw=pLine,text=pMuted,minimum size=.39cm] (k) at (1.28,1.61) {$K$};
  \draw[pDash] (b.south) -- (.40,1.61) -- (k.west);
  \node[pSmall,anchor=west] at (1.57,1.61) {Alternative direction: paused};
  \node[pCard,draw=pTeal,fill=pTeal!5,text width=4.16cm,align=center,
    minimum height=.62cm] (measure) at (3.85,1.00)
    {Measure the complete program on sealed Train};
  \draw[pEvidence] (v.east) -- (6.19,2.96) -- (6.19,1.00) -- (measure.east);
  \node[pSmall,anchor=west,text=pBlue] at (.03,.31) {Best $B$: guarded gains};
  \node[pSmall,anchor=west,text=pTeal] at (3.13,.31) {$W$: evidence + protection};

  \foreach \y/\n/\t/\sub in {
    3.62/1/{Acquired?}/{Inspect tool outputs},
    2.55/2/{Consumed?}/{Inspect routing traces},
    1.48/3/{Interpreted?}/{Inspect criteria and fusion}}{
    \node[pDot,draw=pGold,fill=pGold!8,text=pGold,minimum size=.43cm] at (6.73,\y) {\n};
    \node[pTitle,anchor=west] at (7.09,\y+.07) {\t};
    \node[pSmall,anchor=west] at (7.09,\y-.31) {\sub};
  }
  \draw[pArrow,draw=pGold] (6.73,3.37) -- (6.73,2.81);
  \draw[pArrow,draw=pGold] (6.73,2.30) -- (6.73,1.74);
  \node[pSmall,anchor=west] at (6.48,.40) {Check alternatives.\\Not a success gate.};

  \draw[pRule] (10.12,.11) -- (10.12,4.18);
  \foreach \y/\mark/\t/\sub/\col in {
    3.70/C/Continue/{Evidence; guards hold}/pTeal,
    2.66/P/Pause/{Reallocate search effort}/pBlue,
    1.62/E/Exhausted/{Resource budget ended}/pGold,
    .66/R/Refuted/{Prediction contradicted\\within declared scope}/pRed}{
    \node[pToken,fill=\col,minimum size=.33cm,font=\sffamily\bfseries\fontsize{7}{8}\selectfont]
      at (10.54,\y+.05) {\mark};
    \node[pTitle,text=\col,anchor=west] at (10.84,\y+.10) {\t};
    \node[pSmall,anchor=north west] at (10.84,\y-.13) {\sub};
  }
\end{tikzpicture}

\caption{A direction is a hypothesis, not an edit. (a) Connected tool, routing,
and interpretation changes share a fixed parent. Best and working programs
remain distinct; promotion and continuation both require complete measurement.
(b) Check evidence acquisition, consumption, and interpretation.
(c) Distinguish continuation, pause, exhaustion, and scoped refutation.
The revision branches are schematic, not an observed optimization trajectory.}
\label{fig:directions}
\end{figure}

Search retains both a best development program and a working program for the
active direction (Figure~\ref{fig:directions}a). New directions begin at the best
program. Continuing a direction may retain a fully measured intermediate program
when new evidence supports further revision and the protection conditions hold.
Development promotion requires improved agreement, complete execution,
nondecreasing coverage, and a prespecified stratum-regression guard.
It is a selection rule, not a significance test.

Allocation can request evidence, continue, pause, or switch direction. Resource
exhaustion is distinct from refutation. A direction is refuted only within a
declared scope when evidence contradicts its prediction and relevant alternatives,
including missing observations, unused tools, and incorrect fusion, have been
examined. Failure of one candidate or a nonsignificant interval is insufficient.
On normal completion, an eligible shortlist is frozen from training evidence
before terminal validation. System confirmation uses fresh repeated scoring;
component attribution is a separate controlled comparison. Neither substitutes
for independent Test/OOD evaluation. Appendix~\ref{app:algorithm} gives
Algorithm~\ref{alg:era} and the attribution contract.

\subsection{Pointwise measurement}
\label{sec:measurement}

Each artifact receives $R\geq3$ complete scores under a fixed protocol.
Aggregation precedes comparison:
\begin{equation}
 \bar{s}_\theta(x,y)=\median_{r=1,\ldots,R}s_\theta^{(r)}(x,y),\qquad
 z_i=\ind[\bar{s}_\theta(x_i,y_i^+)>\bar{s}_\theta(x_i,y_i^-)].
 \label{eq:measurement}
\end{equation}
Here $y_i^+$ is the human-preferred artifact, a designation never exposed to the
scorer. Predicted ties are incorrect on strict-preference pairs; human ties
require a separate prespecified policy. This is not majority voting over pairwise
decisions. Missing executions retain their denominators and are reported through
coverage and missing-outcome bounds (Appendix~\ref{app:measurement}).
